\chapter{Nhân Cách Của Người Dạy, Nhân Cách Của Người Học}

\begin{leadbox}
Không có bài học nhân cách nào đủ sức thuyết phục nếu người dạy không sống điều mình giảng.
Và cũng không có sự trưởng thành nào thật sự xảy ra nếu người học chỉ chờ bên ngoài thay đổi. Chương này nói về độ khớp giữa lời nói, hành vi và tâm thế sống.
\end{leadbox}

\storycard{101. Giáo Viên Đi Trễ}{The Late Teacher}{Uy tín không chỉ nằm ở bài giảng hay; nó được dựng lên từ những điều nhỏ được giữ đúng mỗi ngày.}
\storycard{102. Học Sinh Xin Lỗi Thật Lòng}{A Real Apology}{Một lời xin lỗi chỉ có trọng lượng khi nó đi cùng nỗ lực sửa lại điều đã làm lệch cán cân.}
\storycard{103. Nói Về Trung Thực Trong Mùa Điểm}{Talking About Honesty During Grade Season}{Dạy sự trung thực khó nhất khi ai cũng đang bị áp lực thành tích đẩy tới giới hạn.}
\storycard{104. Thầy Nói Ít Nhưng Làm Rõ}{Few Words, Clear Deeds}{Sự mẫu mực âm thầm thường giáo dục mạnh hơn nhiều so với những bài diễn thuyết dài.}
\storycard{105. Trò Tốt Nhưng Kiêu Kiêu}{The Good Student Who Looks Down on Others}{Thành tích cao không tự động đồng nghĩa với một nhân cách đẹp và biết tôn trọng người khác.}
\storycard{106. Một Lời Hứa Với Cả Lớp}{A Promise to the Class}{Khi một lời hứa được giữ đúng, niềm tin của tập thể được bồi đắp nhanh hơn ta tưởng.}
\storycard{107. Giữ Bí Mật Cho Học Sinh}{Keeping a Student's Trust}{Sự tin cậy được hình thành khi người lớn biết giữ ranh giới và bảo vệ phẩm giá của trẻ.}
\storycard{108. Nói Thật Về Sự Thiên Vị}{Admitting Bias}{Nhận ra độ lệch của mình là bước đầu tiên để trở thành một người dạy công bằng hơn.}
\storycard{109. Hành Vi Đẹp Không Có Giám Sát}{Doing Right Unseen}{Nhân cách được nhìn rõ nhất ở những lúc không có ai chấm điểm hay đứng nhìn.}
\storycard{110. Sự Tôn Trọng Hai Chiều}{Respect Both Ways}{Một lớp học mạnh là nơi sự tôn trọng không chỉ đi từ trên xuống, mà còn đi từ dưới lên bằng tự giác.}
